% --- PACOTES BÁSICOS ---
\usepackage{lmodern}            % Usa a fonte Latin Modern
\usepackage[T1]{fontenc}        % Selecao de codigos de fonte.
\usepackage[utf8]{inputenc}     % Codificacao do documento
\usepackage{lastpage}           % Usado pela Ficha catalográfica
\usepackage{listings}           % Para o ambiente lstlisting
%\usepackage{geometry}           %Usado pra alterar as margens pra capa.
\usepackage{indentfirst}        % Indenta o primeiro parágrafo de cada seção.
\usepackage{color}              % Controle das cores
\usepackage{graphicx}
% Inclusão de gráficos e imagens
\usepackage{microtype}          % para melhorias de justificação
\usepackage{mathtools}          % Matemática
\usepackage{multirow}  
\usepackage{longtable}
% Tabelas
%\usepackage{natbib}             % Citações bibliográficas

\usepackage[table,xcdraw]{xcolor} % Tabelas (cores)
\usepackage{float}              % Para que o posicionamento funcione

% --- PACOTES ADICIONAIS ---
\usepackage{epigraph}
\usepackage{amssymb}
\usepackage{subcaption}
\usepackage[export]{adjustbox}
\usepackage{wrapfig}
\usepackage{amsmath}
\usepackage{pdfpages}
\usepackage{rotating}
\usepackage{pdflscape}
\usepackage{lipsum}             % Geração de dummy text
\usepackage{tikz}
\usetikzlibrary{positioning, arrows.meta, calc}
\tikzset{
  mecanismo/.style={
    red,
    line width=0.5pt,
    shorten <=4pt,
    shorten >=4pt,
    -{Stealth[length=2mm, width=1.5mm]}
  }
}
\usetikzlibrary{fit}
\usetikzlibrary{calc}
\usepackage{transparent}
\usepackage{import}
% --- PACOTES DE QUIMICA ---
\usepackage{chemfig}
\setchemfig{atom sep=2em}
\usepackage[version=4]{mhchem}
\usepackage[modules=all]{chemmacros}
\usepackage{float}
\usepackage{tgpagella}
\linespread{1.5}

% --- PACOTES DE CITAÇÕES ---
\usepackage[brazilian,hyperpageref]{backref}  % Paginas com as citações na .bib
\usepackage[alf,abnt-etal-list=0,abnt-etal-cite=6]{abntex2cite} % Citação ABNT